\documentclass[../../main.tex]{subfiles}
\graphicspath{{../../image/}{../../figure/}}
\begin{document}

\section{梯度下降优化}\label{section:ch07-梯度下降优化}

除了最简单的神经网络模型之外，绝大多数网络的误差曲面都具有极其复杂的形态，因此最终求得的解将很大程度上依赖于初始参数值 $\mathbf{w}^{(0)}$ 的选取等因素。为获得一个足够好的解，通常需要多次运行基于梯度的优化算法，每次都选用一个不同的随机初始点，最终通过在一个独立验证集上进行性能评估来确定最优解。

大多数优化技术首先要为权重向量选取一个初始值 $\mathbf{w}^{(0)}$，然后按照如下迭代规则在权重空间中逐步搜索：
\begin{align}
\mathbf{w}^{(\tau)} = \mathbf{w}^{(\tau-1)} + \Delta\mathbf{w}^{(\tau-1)},
\label{align:ch07-iter}
\end{align}
其中 $\tau$ 标记迭代步骤。不同算法对权重向量如何更新 $\Delta\mathbf{w}^{(\tau)}$ 有不同选择。连续非线性函数的优化问题一直是学术界广泛研究的课题。

\subsection{梯度信息的使用}\label{section:ch07-梯度信息的使用}

在~\eqref{align:ch07-taylor}~给出的误差函数二次近似表达式中，误差曲面形态由向量 $\mathbf{b}$ 和矩阵 $\mathbf{H}$ 共同决定，它们总共包含 $W(W+3)/2$ 个独立元素（由于 $\mathbf{H}$ 具有对称性），其中 $W$ 表示权重向量 $\mathbf{w}$ 的维度（即网络中可学习参数总数，见习题~7.7）。因此该二次近似所对应的最小值位置取决于 $O(W^2)$ 个参数，这意味着我们至少需要收集到 $O(W^2)$ 个相互独立的信息片段才能准确定位到该最小值。如果不考虑梯度信息，预计需要 $O(W^2)$ 次误差函数数值评估，而每次评估又需要 $O(W)$ 步计算操作。这样一来，采用这种策略寻找最小值所需的总计算量将高达 $O(W^3)$。

现在将上述情况与利用梯度信息的优化算法进行对比。由于梯度向量 $\nabla E$ 的长度为 $W$，每次计算 $\nabla E$ 都能提供 $W$ 个独立信息片段。据此推断，我们有望仅通过 $O(W)$ 次梯度评估就能找到函数最小值（参见第~8~章）。正如后续章节将展示的那样，利用误差反向传播算法，每次梯度评估的计算复杂度仅为 $O(W)$，因此现在能够在 $O(W^2)$ 的总计算复杂度内找到最小值。尽管二次近似的有效性仅限于最小值附近的局部区域，但这种效率提升具有普遍性。正是基于这一关键优势，利用梯度信息成了当前所有实用神经网络训练算法的基础。

深度网络的误差函数梯度可以通过误差反向传播技术高效计算（参见第~8~章），因此应用这些梯度信息可以显著提高网络训练速度。

\subsection{批量梯度下降}\label{section:ch07-批量梯度下降}

利用梯度信息指导权重更新的一种最朴素方法是在~\eqref{align:ch07-iter}~中按照负梯度方向进行小步长的迭代调整：
\begin{align}
\mathbf{w}^{(\tau)} = \mathbf{w}^{(\tau-1)} - \eta\,\nabla E(\mathbf{w}^{(\tau-1)}),
\label{align:ch07-batchgd}
\end{align}
其中参数 $\eta>0$ 被定义为\textbf{学习率}（learning rate）。在完成每次权重更新后，需要基于新权重向量 $\mathbf{w}^{(\tau)}$ 重新计算梯度，并重复上述过程。在每一次迭代过程中，权重向量都朝着当前位置处误差函数下降最快的方向进行更新，因此这种方法称为\textbf{梯度下降}（gradient descent），或称\textbf{最速下降}（steepest descent）。

注意，由于误差函数是基于整个训练集进行定义的，因此每次计算 $\nabla E$ 都需要遍历所有训练样本。这种一次性处理全部训练数据的技术称为\textbf{批处理}（batch）方法。

\subsection{随机梯度下降}\label{section:ch07-随机梯度下降}

深度学习方法的性能通常会随训练数据集规模扩大而显著提升。然而当训练集中包含海量数据时，批处理方法的效率会急剧下降，因为每一次误差函数或梯度的计算都需要遍历整个数据集。为寻求更高效优化策略，我们注意到：对于一组独立同分布观测样本，基于最大似然估计构建的误差函数可以自然地分解为每个数据点对应项的和：
\begin{align}
E(\mathbf{w}) = \sum_{n=1}^{N} E_n(\mathbf{w}).
\label{align:ch07-sse-decomp}
\end{align}

在处理大规模数据集时，应用最广泛的训练算法是基于梯度下降的\textbf{随机梯度下降}（Stochastic Gradient Descent, SGD）逐次迭代算法~\citep{Bottou2010}。该算法每次仅基于一个数据点来更新权重向量，具体更新规则如下：
\begin{align}
\mathbf{w}^{(\tau)} = \mathbf{w}^{(\tau-1)} - \eta\,\nabla E_n(\mathbf{w}^{(\tau-1)}),
\label{align:ch07-sgd}
\end{align}
其中索引 $n$ 在每次迭代后循环递增。上述更新过程通过循环遍历训练集中数据点来反复执行。对整个训练集进行一次完整遍历称为一个\textbf{训练轮次}（epoch）。当数据以流的形式持续生成时，这种技术有时也称为\textbf{在线梯度下降}（online gradient descent）。

与批量梯度下降相比，随机梯度下降能更高效地处理数据中的冗余。为理解这一点，考虑一个极端例子：复制某数据集的每个数据点，把规模扩大一倍。这种情况下误差函数数值就是原来两倍；如果相应调整学习率的值，则仍可使用原来的误差函数。对于批处理方法，由于数据集加倍，评估批误差函数梯度所需计算量也将加倍，而随机梯度下降算法则完全不受影响。

随机梯度下降的另一个重要特性是其具备逃离局部极小值的能力，因为整个数据集误差函数的驻点通常并不会是每个单独数据点对应误差函数的驻点。

\subsection{小批量方法}\label{section:ch07-小批量方法}

随机梯度下降的一个缺点是：基于单个数据点计算出的梯度可能大大偏离整个数据集上误差函数的梯度。为兼顾计算效率和估计精度，我们折衷一下，每次迭代时使用一小批数据点（称为\textbf{小批量}，mini-batch）来估计梯度。在确定小批量合适大小时需要注意：利用 $N$ 个样本估计均值的误差为 $\sigma/\sqrt{N}$（参见练习~7.8），其中 $\sigma$ 为生成这些数据的分布的标准差。这表明通过增加小批量规模来提升梯度估计精度的收益是递减的——例如将小批量规模增加~100~倍，估计误差仅能降低到原先的十分之一。选择小批量规模时还需要考虑充分利用运行代码的底层硬件架构：例如在某些硬件平台上，小批量规模如果是~2~的幂次（如~64、128、256~等）就可以获得更好性能。

在使用小批量方法时，要注意数据点应当从整个数据集中随机选取。因为在原始数据集中，由于数据采集方式（例如数据点按字母顺序或日期排序），相邻数据点之间可能存在相关性。为消除这种相关性的影响，通常做法是先将整个数据集进行随机打乱，然后按顺序划分为多个小批量；此外在每次完整遍历数据集后，还可以再次打乱数据以确保每个小批量的数据组合都是全新的，这有助于模型跳出局部最优解。即使采用了小批量方法，这种学习算法通常仍被称为"随机梯度下降"。

\subsection{参数初始化}\label{section:ch07-参数初始化}

像梯度下降这样的迭代式优化算法都要求我们为需要学习的参数设定初始值。参数初始化的具体策略会对算法收敛速度以及训练得到网络的泛化能力产生显著影响；遗憾的是目前还没有能够系统性指导参数初始化过程的完善理论。

尽管如此，参数初始化过程中一个关键考虑因素是\textbf{对称性破缺}（symmetry breaking）。考虑网络中一组隐藏单元或输出单元，它们都接收相同输入信号。如果这些单元的所有参数都被赋予相同值（例如全部初始化为零），那么这些参数在训练过程中将始终步调一致地进行更新，并且每个单元都计算出相同函数，导致功能上冗余。为解决这一问题，可以从某个预先定义分布中随机采样来初始化参数，以打破这种对称性。在计算资源允许的情况下，可以采用不同随机初始化策略多次训练网络，最终选择在独立验证集上性能最佳的模型。

用于初始化权重的分布通常要么是范围为 $[-\epsilon,\epsilon]$ 的均匀分布，要么是形式为 $\mathcal{N}(0,\epsilon^2)$ 的零均值高斯分布。$\epsilon$ 值的选择很重要；在各种经验方法中，何恺明初始化~\citep{He2015b}得到了广泛使用。考虑一个网络，其中第 $l$ 层计算以下转换：
\begin{align}
a_i^{(l)} = \sum_{j} w_{ij}^{(l)} z_j^{(l-1)}, \qquad z_i^{(l)} = \mathrm{ReLU}(a_i^{(l)}),
\label{align:ch07-heinit}
\end{align}
其中 $j$ 是发送连接到单元 $i$ 的单元数，ReLU 激活函数由式~\eqref{align:ch06-relu}~给出。假设使用高斯分布 $\mathcal{N}(0,\sigma_w^2)$ 初始化权重，并假设第 $(l-1)$ 层中单元的输出 $z_j^{(l-1)}$ 具有零均值和方差 $s^2$（见习题~7.9），则可以容易地证明
\begin{align}
\mathbb{E}[a_i^{(l)}] = 0, \qquad \mathrm{var}[a_i^{(l)}] = M s^2 \sigma_w^2.
\label{align:ch07-hevar}
\end{align}
其中的因子来自 ReLU 激活函数。理想情况下我们希望确保在从一层传播到下一层时，预激活的方差既不会衰减到零，也不会显著增长。如果因此要求第 $l$ 层单元也具有方差 $s^2$，那么对于一个具有 $M$ 个输入的单元，用于初始化其权重的高斯分布标准差为：
\begin{align}
\sigma_w = \sqrt{\frac{2}{M}}.
\label{align:ch07-hestd}
\end{align}

我们还可以将初始化分布的尺度 $\epsilon$ 视为一个超参数，并在多次训练中探索不同的 $\epsilon$ 值。偏置参数通常设置为小的正值，以确保大多数预激活在学习开始时处于激活状态——这对 ReLU 特别有帮助，我们希望预激活为正值，以便有非零梯度来驱动学习过程。

除上述方法外，另一类重要参数初始化策略是利用在其他相关任务上预训练得到的网络权重，或者利用各种无监督学习方法。这些技术都属于迁移学习（transfer learning）的范畴（参见~6.3.4~节）。

\end{document}
